\documentclass[11pt]{article}

% ---------------------------------------------------------------------------
% The record of the sweden-residents line: design, decisions, rules and
% outcomes of the cohort-exposure design. Replaces
% output/documents/empirical-strategy/empirical_strategy.tex (2026-11-18).
% cohort-exposure-slides.tex summarises it and follows it.
% Compiles on its own (Overleaf): references in references.bib.
% Metadata for the repository's docs index: prespec-cohort-exposure.md.
% ---------------------------------------------------------------------------

\usepackage[margin=1in]{geometry}
\usepackage{amsmath}
\usepackage{amssymb}
\usepackage{array}
\usepackage{booktabs}
\usepackage{enumitem}
\setlist{itemsep=0.15em, topsep=0.3em}
\usepackage[colorlinks=true, linkcolor=blue, citecolor=blue, urlcolor=blue]{hyperref}
\usepackage[style=authoryear, backend=biber, maxcitenames=2, uniquename=false, uniquelist=false]{biblatex}
\addbibresource{references.bib}

\setlength{\bibitemsep}{0.3em}
\renewcommand*{\bibfont}{\small}
\setcounter{secnumdepth}{2}
\newcolumntype{L}[1]{>{\raggedright\arraybackslash}p{#1}}
\newcommand{\prule}[1]{\item[\(\Diamond\) #1]}

\title{Growing up behind a noise barrier:\\ the cohort-exposure design (pre-specification)}
\author{}
\date{Draft, 19 November 2026}

\begin{document}
\maketitle

\noindent\textbf{Status.} Draft. Nothing has been estimated. Part~\ref{sec:prespec} lists the
open decisions (\S\ref{sec:decisions}) and the rules (\S\ref{sec:rules}); no outcome is loaded
before the decisions are taken.

\noindent\textbf{How to use this document.}
\begin{itemize}
\item It is the line's record: the design, the decisions, the rules fixed before estimation,
  and the outcomes once they come back.
\item A rule changes only through the change log (\S\ref{sec:changelog}), with a date and a
  reason, never after seeing a result the change could be fitted to.
\item The slides follow this document; change the design here first.
\item \textbf{Where it runs.} The 100\,m treatment grid and the stack file (\S\ref{sec:data})
  are built in the project repository. The register linkage and all estimation run with the
  register data (MONA), outside the repository.
\end{itemize}

\tableofcontents

% ===========================================================================
\section{Question and contribution}

\begin{itemize}
\item \textbf{Question.} Does living with road or rail noise during childhood lower grade-9
  scores, and at which ages does it matter?
\item \textbf{Why noise barriers.} Trafikverket has built hundreds of road and rail barriers
  over three decades. Each cuts noise sharply on one side of the road, in a known year, and
  leaves the other side as it was.
\item \textbf{Why Swedish registers.} Every child has a grade-9 score and is linked to their
  mother, whose residence is recorded each year on a 100\,m grid. We can place a family before
  the barrier and follow the child after the family moves.
\item \textbf{Contribution.} A cohort design for noise and school achievement: exposure
  assigned before the barrier, compared within each corridor.
\item \textbf{Relation to the school-level line.} That line measures exposure at school. This
  one measures exposure at \emph{home}. Where a child's home and school are both near barriers,
  the two exposures can enter together.
\end{itemize}

% ===========================================================================
\section{Setting}

\begin{itemize}
\item The Swedish road administration (SRA) provides two kinds of noise mitigation: facade
  insulation, the more common one, and noise barriers.
\item Mitigation is rationed by noise level, using thresholds. It is also provided regardless
  of noise level to groups of properties in one area, or together with road construction
  projects.
\item Residents can apply and can opt out. Some eligible dwellings stay untreated, and
  dwellings worth less than the installation cost are not treated
  \autocite{lindgrenSoundEnvironmentHealth2021}.
\end{itemize}

Two features shape the design:
\begin{enumerate}
\item \textbf{A barrier is a bundled treatment.} It usually arrives with a road project.
\item \textbf{``Untreated'' is inferred from a missing record.} Untreated dwellings may have
  private insulation or SRA insulation.
\end{enumerate}

% ===========================================================================
\section{What we take from the literature}\label{sec:literature}

\begin{center}
\footnotesize
\renewcommand{\arraystretch}{1.25}
\begin{tabular}{@{}L{0.31\textwidth}L{0.3\textwidth}L{0.32\textwidth}@{}}
\toprule
\textbf{Study} & \textbf{We take} & \textbf{We do differently} \\
\midrule
\textcite{lindgrenSoundEnvironmentHealth2021}: Swedish noise mitigation and health & The setting & Keep movers; same-road controls, not never-treated \\
\textcite{morettiTrafficNoiseExternality2025}: barriers and house prices & The 600--1,500\,m band & Follow people, not dwellings; no unbuilt barriers in Sweden \\
\textcite{greenhillNoisePollutionInfant2026}: noise and infant health & Noise and air pollution move together & Rail as a near noise-only treatment \\
\textcite{gronqvistUnderstandingHowLow2020}: lead, Swedish cohorts & Cohort \(\times\) exposure, threshold fixed in advance, percentile outcome & Place-specific timing; dose from years lived \\
\textcite{hollingsworthThousandCutsCumulative2025}; \textcite{berkhoutLeadExposureCognitive2025} & Duration against intensity; age profile; distance cut-off from the data & Noise, not lead \\
\bottomrule
\end{tabular}
\end{center}

\paragraph{\textcite{lindgrenSoundEnvironmentHealth2021}} work in almost our setting:
treated and untreated residents of single-family homes along state highways the SRA
considered for mitigation (apartments dropped). They keep individuals born before 1975 who
moved into a not-yet-mitigated dwelling and stayed at the same address throughout 2003--2017
(10,113 kept, 31,379 dropped for short residence). Controls are considered dwellings not
mitigated by 2017. Three lessons:
\begin{enumerate}
\item \textbf{Do not condition on staying.} The stayers have a mean residence of 19.5 years
  against 5.5 among those dropped. If mitigation affects moving, this conditions on an outcome
  of treatment. We assign treatment by residence \emph{before} the barrier and keep everyone
  who moves later (intention to treat).
\item \textbf{Test for pre-trends.} Their main sample has flat pre-trends, but disease risk in
  the barrier group already falls before treatment (their Figure 4b). Our design has a natural
  placebo: cohorts tested before the barrier.
\item \textbf{Do not rely on never-treated controls alone.} Eligible dwellings that are never
  treated may have a persistently low priority correlated with their trends. We use the same
  road at the same distance, and not-yet-treated corridors as a secondary comparison.
\end{enumerate}

\paragraph{Insulation in the comparison group.} Untreated cells may hold private or SRA
insulation. Restricting treatment to barriers does not remove it.
\begin{itemize}
\item If private take-up is highest closest to the source, contamination grows with exposure:
  bias toward zero.
\item If private retrofits are mostly cosmetic or for energy in quieter dwellings, the sign is
  unclear.
\item Owner-occupiers retrofit more than renters (split incentives).
\item We discuss the likely direction and use insulation records where they exist.
\end{itemize}

\paragraph{\textcite{morettiTrafficNoiseExternality2025}} study barriers alone. They compare
dwellings 0--500\,m behind a barrier with dwellings 500--1,500\,m behind it, and add a third
difference against barriers proposed but never built. Two features do not carry over:
\begin{itemize}
\item Repeated house sales: which houses sell may respond to treatment.
\item Treatment and outcome are tied to dwellings, not to the people who may move, so the
  estimate mixes the direct effect with sorting.
\end{itemize}
They test baseline differences between built and unbuilt projects but not why projects were
dropped, so differential trends remain possible. Trafikverket records only built barriers, so
the unbuilt comparison is not available in Sweden.

\paragraph{\textcite{greenhillNoisePollutionInfant2026}} stress confounding with air pollution.
Field studies find lower near-road pollution behind roadside barriers
\autocite{baldaufImpactsNoiseBarriers2008}, so a road barrier does not change only noise. We
(i) read the road-barrier estimate as the effect of the bundle and (ii) use rail barriers as a
near noise-only treatment, since most Swedish rail traffic is electric.

\paragraph{The lead literature} uses the same cohort-exposure logic.
\textcite{gronqvistUnderstandingHowLow2020} interact birth cohort with the parish's air lead
level before the 1980--81 cut in gasoline lead. \textcite{hollingsworthThousandCutsCumulative2025}
measure lifetime lead from nearby racetracks, scaled by distance, around the 2007 switch to
unleaded racing fuel. \textcite{berkhoutLeadExposureCognitive2025} combine age at first exposure
to a contaminated site with distance to it. Four lessons:
\begin{enumerate}
\item \textbf{Exposure varies in intensity, not only in timing.} A dose-response also separates
  how long from how much. We add each cell's predicted noise reduction (\S\ref{sec:intensity}).
\item \textbf{A threshold fixed in advance turns the event study into a test.} Grönqvist et al.\
  locate a threshold and predict effects only above it. We do the same with the noise level
  before the barrier.
\item \textbf{The event study traces an age profile.} Each later cohort adds one year of
  exposure at a younger age. The earliest ages matter most in these studies, in utero most of
  all \autocite{berkhoutLeadExposureCognitive2025}.
\item \textbf{Distance cut-offs come from the data.} Hollingsworth et al.\ set their radius where
  a distance spline fades; plant emissions affect air quality only within a mile
  \autocite{currieEnvironmentalHealthRisks2015}. We estimate the profile behind the barrier
  (\S\ref{sec:comparison}).
\end{enumerate}

% ===========================================================================
\section{Design}\label{sec:design}

A \textbf{cohort exposure design} \autocite{dufloSchoolingLaborMarket2001}, as a
difference-in-differences within each barrier:
\begin{itemize}
\item \textbf{Across sides:} children from protected cells against children from comparison
  cells along the same road (\S\ref{sec:comparison}).
\item \textbf{Across cohorts:} cohorts tested before the barrier (zero exposure, the placebo)
  against later cohorts. Exposure rises one year per cohort for cohorts born before the
  barrier, and is complete, including the whole pregnancy, from the cohort born a year after
  it.
\item \textbf{Assignment before the project:} each child is placed where the family lived
  before the barrier and its road project were planned, and followed afterwards.
\end{itemize}

\subsection{Notation}

\begin{itemize}
\item \(i\): child; \(f(i)\): the child's household (mother); \(a_i\): birth year.
\item \(A = 16\): age at the grade-9 test for a child who started school on time and never
  repeated or skipped a grade. \(t^*_i = a_i + A\): the intended test year.
\item \(b\): a barrier, built in year \(\tau_b\). \(\rho_b\): its reference year.
\item \(g_0(i)\): the reference cell, the 100\,m grid cell where the mother lived on
  31~December of \(\rho_b\). \(P_i \in \{0,1\}\): 1 if that cell is protected by \(b\) (behind
  the barrier, on the side it shields: the \texttt{protected} tier).
\item \(e_i = t^*_i - \tau_b\): event time. \(D_i = \min\{\max(0, e_i),\, A\}\): the dose,
  years lived with the barrier before the test.
\item Reference cohort: \(e = -2\). Cohorts at \(e = -1\) and \(e = 0\) are transition cohorts
  (\S\ref{sec:predictions}).
\item \(r(i)\): the reference cell's distance band from the road.
\item \(K\), \(L\): the window, test years \([\tau_b - K,\ \tau_b + A + L]\).
\end{itemize}

\subsection{Assignment}\label{sec:assignment}

\paragraph{The problem.} Families may move in, move out or stay \emph{because of} the barrier.
Projects are announced years ahead, so the current address would mix sorting into the ramp.

\paragraph{The rule.}
\begin{itemize}
\item Each child gets the mother's cell at \(\rho_b\) and keeps it whatever the family does
  later. Assignment is predetermined; later moves are part of the intention-to-treat effect,
  not selection.
\item Children born after \(\tau_b\) have no residence of their own before the barrier, yet
  they identify full exposure. Assigning through the mother gives them a cell, like their
  older siblings.
\item \(\rho_b\) balances anticipation against dilution:
  \begin{itemize}
  \item too late (e.g.\ \(\tau_b - 1\)) admits families who already reacted to the
    announcement or the noisy construction work;
  \item too early (e.g.\ the child's birth) lets unrelated moves dilute assignment: more
    assigned children never live behind the wall, a weaker first stage, an attenuated estimate.
  \end{itemize}
\item \(\rho_b\) is the year before the road project started or was announced, whichever is
  earlier, from project documents; \(\tau_b - 3\) where no record exists. No historical source
  has been found (Trafikverket's traffic-information service keeps finished roadworks for about
  a year), so \(\tau_b - 3\) is the rule for every barrier until one is.
\item \textbf{Robustness:} \(\rho_b = \tau_b - 1\) and \(\tau_b - 5\). They are compared per
  year lived (2SLS, \S\ref{sec:mobility}), which dilution does not move. If anticipation
  matters, later reference dates give smaller estimates.
\item \textbf{Anticipation} is measured directly: moves out of and into the reference cell
  between \(\rho_b\) and \(\tau_b\), by side.
\end{itemize}

\paragraph{Intended, not actual, test year.} Event time and dose use \(t^*_i\). A late school
start or a repeated grade may itself respond to the barrier, so the actual test year would make
them partly an outcome. Off-age children stay in their birth cohort and are scored when they
take the test; the gap of about a year is part of the ITT. Grade 9 on time is a second outcome,
and we re-estimate on children tested in their intended year.

\subsection{Comparison groups}\label{sec:comparison}

\begin{enumerate}
\item \textbf{Opposite side, same distance (main).} Cells at the same distance from the same
  road segment, on the side the barrier does not shield (sides from the barrier audit). They
  share the baseline noise level, the road project, traffic changes and the local economy.
  \begin{itemize}
  \item The opposite side may lack a barrier \emph{because} it has few dwellings or was
    insulated; we check this.
  \item Of barriers built by 2007, the opposite side is more than half protected for 37\% of
    road and 42\% of rail barriers, nearly always by the same project (within 2 years).
  \item Weighted by children it is thinner: stacks with at least one
    expected child per cohort across the road hold 27\% (road) and 45\% (rail) of protected
    children (\S\ref{sec:power}).
  \end{itemize}
\item \textbf{Farther behind the barrier (600--1,500\,m band)}, after
  \textcite{morettiTrafficNoiseExternality2025}. Quieter at baseline and less affected by
  traffic changes. It keeps two thirds of protected children. Pooled with group~1 in the main
  estimate if decision~\ref{dec:comparison} is taken as recommended; otherwise secondary.
\item \textbf{Not-yet-treated corridors}, with a later \(\tau\): robustness on group 1. We do
  not rely on never-treated places alone: places that qualified but were never shielded may have
  low priority for reasons that also shape their trends.
\end{enumerate}

\paragraph{Distance bands} \(r\): 0--100, 100--250 and 250--600\,m, fixed in advance.

\paragraph{Distance profile.} In the parametric specification (\S\ref{sec:dose}), \(P_i\) is
replaced with 50\,m bins of the reference cell's distance from the wall, up to 1,500\,m, each
against the opposite side at the same distance.
\begin{itemize}
\item It shows how far protection reaches, and whether the 600--1,500\,m band is unaffected,
  which group 2 requires (decision~\ref{dec:profile}).
\item For cohorts tested before the barrier it is a balance check: scores should vary with
  distance in the same way on both sides \autocite[as in][]{berkhoutLeadExposureCognitive2025}.
\end{itemize}

\subsection{Stacks}\label{sec:stacks}

Staggered construction makes the comparison the central choice. A staggered estimator drawing
on the national pool of never- or not-yet-treated units \autocite[as in the school-level
analysis,][]{callawayDifferenceinDifferencesMultipleTime2021} would compare a barrier with cells
near other roads, without netting out the corridor's own project, traffic and economy. We
therefore estimate a \textbf{stacked event study} \autocite{cengizEffectMinimumWages2019}:
\begin{itemize}
\item \textbf{One stack per barrier}: children whose mother's reference cell lies in \(b\)'s
  corridor (protected and comparison cells) and whose test year falls in the window. Children
  stay in whether or not the family moves.
\item \textbf{Clean controls}: comparison cells protected by another barrier inside the window
  are dropped; no already-treated cell is a control.
\item \textbf{Weights.} OLS on a stack weights barriers by size and balance, not by a chosen
  estimand. We use the corrective weights of \textcite{wingStackedDifferenceinDifferences2024}:
  protected children weight 1, comparison children in stack \(b\)
  \[ Q_b = \frac{N^P_b / N^P}{N^C_b / N^C}, \]
  with \(N^P_b, N^C_b\) the protected and comparison children in the stack and \(N^P, N^C\)
  their sums over stacks. The estimate is an ATT per protected child: a barrier shielding a
  dense neighbourhood counts more. Stacks are successive cohorts, not a balanced panel, so the
  weights are computed within each barrier \(\times\) event-time cell (an adaptation of their
  formula). Equal barrier weights: robustness.
\item \textbf{Two balanced panels.} Each \(\beta_e\) averages the barriers observed at event
  time \(e\). If barriers enter or leave across \(e\), the shape partly reflects which barriers
  are in the sample, since old and recent walls differ in height, road type and region. Like
  Wing et al., who trim to stacks observed over the whole window, we balance, twice:
  \begin{itemize}
  \item \textbf{Short window, \(e = -K, \dots, +3\)}, for the step (\S\ref{sec:predictions}),
    the main estimate. Every dated barrier built at least three years before the last grade
    year enters, including the many recent walls (71\% of dated road and 54\% of dated rail
    records were built 2015 or later).
  \item \textbf{Full window, \(e = -K, \dots, A + L\)}, for the ramp and the in-utero step. It
    needs scores up to \(\tau_b + A + L\), at least 17 years after construction. With grade-9
    data for 1996--2022 and the recommended \(K = 4\), \(L = 1\) (decision 1), that is
    \textbf{barriers built 2000--2005}. A plateau (\(L \ge 2\)) is out of reach for almost all
    walls.
  \end{itemize}
  The event study on all dated barriers, unbalanced, is a robustness check.
\item \textbf{Clusters.} Corridors overlap, so a comparison child can enter several stacks.
  Standard errors are clustered on groups of connected corridors (barriers whose corridors share
  a cell); with a moderate number of groups we also report wild cluster bootstrap and
  randomisation-inference p-values.
\end{itemize}

\subsection{Sample size and power}\label{sec:power}

Planning numbers from open data, computed in the repository on 19 November 2026 (notebook
\texttt{21\_population\_power}). No outcome is loaded. The register counts of rule G4 replace
them.
\begin{itemize}
\item \textbf{Children per cohort} in a cell at \(\rho_b\): GHS-POP residents on the 100\,m grid
  (modelled, interpolated between epochs) times the cell's share aged 0--19 in SCB's 1\,km grid,
  divided by 20. Children assigned through the mother's cell are approximated by the children
  living there.
\item \textbf{Usable stack}: at least one expected child per cohort on each side.
\item \textbf{MDE}: 5\% two-sided, 80\% power, outcome in SD, no covariates, \(K = 4\), the
  weights of \S\ref{sec:stacks}. Comparison cells shared by several stacks are treated as
  independent, so the numbers are optimistic. A within-cell correlation of 0.02 raises them by
  about 5\%; covariates explaining a quarter of the variance lower them by about 13\%.
\end{itemize}

\begin{center}
\footnotesize
\begin{tabular}{@{}llrccc@{}}
\toprule
& & protected children & \multicolumn{3}{c}{share kept / MDE (SD)} \\
\cmidrule(l){4-6}
window & kind & per cohort & opposite side & band & both \\
\midrule
short, built 2000--2019 & road & 523 & 27\% / 0.46 & 67\% / 0.28 & 78\% / 0.22 \\
                        & rail & 1,186 & 45\% / 0.25 & 66\% / 0.18 & 78\% / 0.15 \\
full, built 2000--2005  & road & 16 & \multicolumn{3}{c}{4 usable stacks: not estimable} \\
                        & rail & 410 & 42\% / 0.38 & 59\% / 0.33 & 75\% / 0.26 \\
\bottomrule
\end{tabular}

\smallskip
\emph{Short window: \(\delta^{\text{SR}}\); full window: the step \(\delta\). The rail full window
detects \(\beta\) per year of dose of 0.019 / 0.017 / 0.013 SD.}
\end{center}

\paragraph{If every barrier had a construction year.} Half the road and a quarter of the rail
records are undated (\S\ref{sec:data}), and an undated barrier has no stack. The registers bound
the missing years: 978 of the 1,114 undated road records were in NVDB's first load in 2006, and
451 of the 487 undated rail records carry BIS's first-load connection date (before April 2019).
\emph{Expected}: each missing year drawn from dated barriers of the same register vintage, within
the bound. \emph{Ceiling}: every undated barrier inside the window, within the bound. MDE in SD,
opposite side / band / both:

\begin{center}
\footnotesize
\begin{tabular}{@{}lccc@{}}
\toprule
& today & expected & ceiling \\
\midrule
road, short-run \(\delta^{\text{SR}}\) & 0.46 / 0.28 / 0.22 & 0.33 / 0.22 / 0.17 & 0.25 / 0.17 / 0.14 \\
rail, short-run \(\delta^{\text{SR}}\) & 0.25 / 0.18 / 0.15 & 0.17 / 0.12 / 0.10 & 0.16 / 0.12 / 0.10 \\
road, full-window \(\delta\) & not estimable & 0.53 / 0.36 / 0.31 & 0.29 / 0.20 / 0.17 \\
rail, full-window \(\delta\) & 0.38 / 0.33 / 0.26 & 0.28 / 0.20 / 0.16 & 0.21 / 0.14 / 0.12 \\
\bottomrule
\end{tabular}
\end{center}

\paragraph{What follows.}
\begin{itemize}
\item As drafted (opposite side, road and rail apart), the design detects only large effects:
  about 0.25 SD for rail and 0.46 SD for road in the short run.
\item Pooling the band with the opposite side brings this to 0.15 and 0.22 SD
  (decision~\ref{dec:comparison}). Pooling road and rail as well gives 0.13 SD, little beyond
  rail alone, and mixes two treatments (decision~\ref{dec:primary}).
\item The missing construction years are the largest gain. The rail years add most to the short
  window: with them, rail's protected children in usable stacks roughly double. The road years
  of the 2006 load add most to the full window, which they open to road walls; 57\% of that
  load's dated walls were built before 2000, outside every window.
\end{itemize}

% ===========================================================================
\section{Specification}\label{sec:spec}

\subsection{Main: stacked cohort event study}

\[
S_i = \sum_{e \neq -2} \beta_e \,\mathbf{1}\{e_i = e\}\, P_i \;+\; \gamma_{b, P_i, r(i)} \;+\; \eta_{b, e_i} \;+\; X_{f(i)}'\phi \;+\; \varepsilon_i
\]

\begin{itemize}
\item The reference cohort is \(e = -2\), not \(e = -1\): cohorts at \(e = -1\) and \(0\) may
  already be affected by the construction works (\S\ref{sec:predictions}).
\item \(S_i\): the child's \textbf{percentile rank on the grade-9 merit value} among all children
  tested in the same year \autocite[as in][]{gronqvistUnderstandingHowLow2020}. Grading changed
  in our period (goal-referenced grades from 1998, then the A--F scale and a higher maximum
  merit value in the 2010s; exact years to be confirmed with SCB), and a window spans 20 or more
  test years. Barrier \(\times\) year FE absorb a shift in level but not in spread; ranking
  handles both. Secondary: the raw merit value and national test results, ranked the same way;
  grade 9 on time.
\item \(\gamma_{b,P,r}\): barrier \(\times\) side \(\times\) distance band FE. Permanent level
  differences between sides and bands within a corridor.
\item \(\eta_{b,e}\): barrier \(\times\) event time FE, the same as barrier \(\times\) test year
  within a stack. They absorb every corridor-wide shock in a year: traffic, local shocks, and the
  part of the project common to all distances. They pool bands, so a project effect that fades
  with distance stays in \(\beta_e\) if the two sides spread differently across bands.
  \textbf{Robustness:} barrier \(\times\) band \(\times\) event time FE \(\eta_{b,r,e}\), so
  every comparison is across the road at the same distance in the same year; bands with
  children on one side only drop out (we report the protected children lost), and the weights
  are computed within barrier \(\times\) band \(\times\) event time. The 600--1,500\,m band has
  no protected children, so this is the opposite-side-only estimate; with the band pooled in
  the main estimate (decision~\ref{dec:comparison}), rule B3 compares the two.
\item \(X_{f(i)}\): parental education, log family income and immigrant background at
  \(\rho_b\).
\end{itemize}

\paragraph{No school fixed effects.}
\begin{itemize}
\item The grade-9 school is chosen after \(\tau_b\). Actual-school FE would remove changes of
  school the barrier causes, and would bring back the sorting that assignment at \(\rho_b\)
  keeps out.
\item The cost: teachers set the merit value, leniency differs across schools and over time,
  and roads often split catchments. A grading change at one side's school is not absorbed by
  \(\eta_{b,e}\) and loads onto \(\beta_e\).
\item \textbf{Robustness:} FE for the \emph{assigned} school \(\times\) test year, fixed before
  treatment: the most common grade-9 school of children from the reference cell tested before
  \(\rho_b\), or an older sibling's school.
\item We report how many protected and comparison children in a corridor attended the same
  grade-9 school before \(\tau_b\). Shared schools: the concern is small. Separate schools:
  actual-school FE would leave only families who cross the road for school.
\item \textbf{Decomposition, not robustness:} actual school \(\times\) test year FE. The gap to
  the ITT shows how much runs through changes of school and peers, next to the event study with
  a change of school as the outcome.
\end{itemize}

\subsection{What the event study should show}\label{sec:predictions}

\textbf{Our prior: noise acts at once.} The wall lowers noise at home immediately, and with it
sleep loss and disturbed concentration and homework. So we expect a \textbf{step}:
\begin{itemize}
\item \(\beta_e = 0\) for \(e \le -2\): tested before the barrier, the placebo (reference
  \(e = -2\)).
\item \textbf{Transition, \(e = -1\) and \(e = 0\).} Construction works are louder than the road
  was. If noise acts fast, they may lower the scores of cohorts tested during the works (a dip).
  The \(e = 0\) cohort is tested in the construction year, and the wall may be finished only
  after the spring term, so it is partly treated at most. Both are estimated on their own, in
  neither the placebo nor the step.
\item \textbf{Step at \(e = 1\).} From \(e = 1\) on, every cohort spends the whole grade-9 year
  in a quieter home, so an immediate effect raises all of them by the same \(\delta\). The
  merit value reflects learning over the last school years, not only the test year, so the step
  may build up over \(e = 1, \dots, 3\) rather than appear at once.
\item \textbf{Ramp on top}, if exposure also accumulates: \(\beta_e\) rises for
  \(1 \le e \le A\). Cohort \(e\) was exposed from age \(A - e\), so
  \(\beta_e - \beta_{e-1}\) is one more year of exposure at age \(A - e\). The ramp is an age
  profile: linear if every year matters equally, steepest at high \(e\) if early childhood
  matters most.
\item \(\beta_{A+1} - \beta_A\): exposure in utero. The cohort at \(e = A\) was born in the
  construction year; from \(e = A + 1\) on, exposed throughout pregnancy.
\item \(\beta_e\) flat for \(e \ge A + 1\): the same dose. Later, adaptation, traffic growth or a
  wall on the other side may bend it. With grade data to 2022 few walls reach beyond
  \(e = A + 1\) (decision 1), so the plateau is described, not tested.
\end{itemize}

Zero, step, then a ramp or a plateau is hard to produce through confounding. The shapes read as
follows:

\begin{center}
\small
\begin{tabular}{@{}ll@{}}
\toprule
\textbf{Shape} & \textbf{Points to} \\
\midrule
Step at \(e = 1\), flat after & an immediate effect of noise (our prior) \\
Ramp from \(e = 1\), no step & cumulative exposure only \\
Pre-trend before \(e = -1\) & sorting, anticipation \\
Dip at \(e = -1, 0\) only & construction noise \\
Ramp, no plateau after \(e = A + 1\) & trends set off by the road project \\
\bottomrule
\end{tabular}
\end{center}

\subsection{Effect per year of dose}\label{sec:dose}

\paragraph{Short window: the main estimate.} On the short panel (\(e = -K, \dots, +3\)):
\begin{align*}
S_i ={}& \delta^{\text{SR}}\, \mathbf{1}\{1 \le e_i \le 3\}\, P_i + \sum_{e \in \{-1, 0\}} \beta_e\, \mathbf{1}\{e_i = e\}\, P_i + \beta^{\text{pre}}\, \mathbf{1}\{e_i \le -2\}\,(e_i + 2)\, P_i \\
&+ \gamma_{b,P_i,r(i)} + \eta_{b,e_i} + X_{f(i)}'\phi + \varepsilon_i
\end{align*}
\(\delta^{\text{SR}}\) is the short-run effect: the step plus at most three years of any ramp.
\(\beta_{-1}, \beta_0\) are the transition cohorts; \(\beta^{\text{pre}}\) tests for a pre-trend.

\paragraph{Full window: step and ramp.} On the full panel, with flat pre-trends and a roughly
linear ramp, the event study is summarised by
\begin{align*}
S_i ={}& \delta\, \mathbf{1}\{e_i \ge 1\}\, P_i + \beta\, D_i P_i + \sum_{e \in \{-1, 0\}} \beta_e\, \mathbf{1}\{e_i = e\}\, P_i + \beta^{\text{utero}}\, \mathbf{1}\{e_i > A\}\, P_i \\
&+ \beta^{\text{pre}}\, \mathbf{1}\{e_i \le -2\}\,(e_i + 2)\, P_i + \gamma_{b,P_i,r(i)} + \eta_{b,e_i} + X_{f(i)}'\phi + \varepsilon_i
\end{align*}
\begin{itemize}
\item \(\delta\): the step, the same for every cohort tested after the barrier;
  \(\beta\): the effect of each further year of assigned dose; \(\beta^{\text{utero}}\): added
  effect of exposure in pregnancy; \(\beta^{\text{pre}}\): pre-trend test.
\item \(\delta\) and \(\beta\) are separately identified: \(\delta\) from the jump at \(e = 1\),
  \(\beta\) from the slope across the many cohorts after it.
\item \textbf{Age profile:} the slope may differ by the age of the added year, with kinks at 3,
  7 and 13 (bins 0--2, 3--6 before school, 7--12, 13--15).
\item \textbf{Why not regress on \(D_i\) alone?} Cohorts tested before would pool with
  comparison cells at \(D_i = 0\), losing the placebo; and a continuous dose with staggered
  timing mixes comparisons across dose levels and slopes, needing stronger parallel trends
  \autocite{callawayDifferenceinDifferencesContinuousTreatment2024}.
\item \(D_i\) is always the dose assigned from the reference cell. Years actually lived behind a
  barrier reflect moves the barrier may cause; they enter only the estimate per realised year
  (\S\ref{sec:mobility}).
\end{itemize}

\paragraph{Interpretation.} \(\beta_e\) is an intention-to-treat effect of living behind a
barrier as of \(\rho_b\). It includes effects through the family's later decisions: moving,
changing school, changing housing.

\paragraph{Other estimators.} Callaway--Sant'Anna on the stacked data with corridor-specific
controls (should reproduce the estimate); national Callaway--Sant'Anna with not-yet-treated
cells (comparison group 3, comparable with the school-level line);
\textcite{sunEstimatingDynamicTreatment2021}; \textcite{borusyakRevisitingEventStudyDesigns2024}.

% ===========================================================================
\section{Noise or the road project?}\label{sec:noise}

Barriers come with a road project and may lower air pollution
\autocite{baldaufImpactsNoiseBarriers2008,greenhillNoisePollutionInfant2026}.
\begin{itemize}
\item Construction and traffic changes hit both sides, so the opposite side nets them out.
\item Air pollution partly passes over a wall, so the gap between sides is smaller than for
  noise.
\item What remains are the side-specific parts of the project and of air pollution. Three tests
  address them:
\end{itemize}

\begin{center}
\footnotesize
\renewcommand{\arraystretch}{1.2}
\begin{tabular}{@{}L{0.3\textwidth}L{0.13\textwidth}L{0.24\textwidth}L{0.24\textwidth}@{}}
\toprule
\textbf{Test} & \textbf{Rules out} & \textbf{If it is noise} & \textbf{If it is not} \\
\midrule
\textbf{Rail vs.\ road} & Air pollution & Similar ramps (trains mostly electric) & Ramp for road barriers only \\
\textbf{Threshold}, fixed in advance: \(L^0\) above the WHO guideline (53\,dB road, 54\,dB rail) & Project & Ramp above, none below & Ramp also below, in quiet cells as close to the road \\
\textbf{Intensity}: predicted dB removed at each cell & Project & Larger ramp where more dB are removed & Same ramp whatever the dB cut \\
\bottomrule
\end{tabular}
\end{center}
Comparing cells at the same distance holds the project fixed. Air pollution moves with distance
and traffic like noise; only rail against road tests it.

\subsection{Exposure intensity}\label{sec:intensity}

\(P_i\) treats every protected cell alike, although a tall wall by a busy road cuts noise far
more than a low wall 150\,m away. For each protected cell \(g\) we predict
\begin{itemize}
\item \(L^0_g\): the noise level before the barrier, and
\item \(\Delta L_g\): the reduction the barrier causes, both in dB,
\end{itemize}
from inputs fixed before the barrier only: wall height and length, the cell's distance to the
road and the wall, traffic before construction (ÅDT or train counts), with a standard method
such as the Nordic prediction method. So \(\Delta L_g\) does not reflect where dwellings are or
who lives there \autocite[the logic of the predicted lead change in][]{gronqvistUnderstandingHowLow2020}.
Comparison cells get \(\Delta L = 0\). Four uses:
\begin{enumerate}
\item \textbf{Intensity event study.} \(P_i\) replaced by \(\Delta L_{g_0(i)}\): \(\beta_e\) per dB.
  The project reaches every protected cell but is not scaled by \(\Delta L\), so a dose-response
  is hard to produce through it, but only at equal distance, since project effects (construction,
  traffic, air) fade with distance like \(\Delta L\). Hence \(\eta_{b,r,e}\). Variation within a
  barrier: position along the wall (ends gain less), changes in height; across barriers: height
  and traffic. Taller walls go with busier roads and bigger projects, so event time is also
  interacted with pre-construction traffic.
\item \textbf{Duration against intensity.} With dB-years \(Q_i = D_i\, \Delta L_{g_0(i)}\),
  indicators for years of exposure controlling for \(Q_i\)
  \autocite{hollingsworthThousandCutsCumulative2025}: do many years of a small cut matter as much
  as a few of a large one? Identified because barriers at the same event time differ in
  \(\Delta L\).
\item \textbf{Threshold.} Protected cells split at the WHO guideline, 53\,dB \(L_{\text{den}}\)
  road and 54\,dB rail, fixed in advance. A ramp above and none below; with \(\eta_{b,r,e}\),
  cells below get as much of the project but were quiet, so a ramp there points to the project
  or composition. Caveat: at a given distance \(L^0\) differs only with traffic, so the split is
  mostly busier against quieter roads, and the event time \(\times\) traffic interaction would
  absorb it; the threshold guards less well than intensity against project effects that grow
  with traffic. Its value comes from fixing the split before looking.
\item \textbf{Permutation.} \(\Delta L\) permuted across protected cells within each corridor:
  if the effect comes from the dose, the permuted dose-response is flatter.
\end{enumerate}

\paragraph{Feasibility} (register checked 2026-11-02). The prediction is validated against EU
strategic noise maps where they cover a corridor before and after construction. It needs wall
height and distance from the road:
\begin{itemize}
\item \textbf{Rail:} height for 79\% of records (323 of the 414 built by 2007); distance from the
  track centre for 44\%.
\item \textbf{Road:} height missing for 62\% of records and for 138 of the 140 built by 2007;
  walls are stored on the road centreline, so distance is unknown except for walls matched to
  OpenStreetMap.
\item So intensity tests are feasible for rail, not for road unless heights come from elsewhere.
  Without heights \(\Delta L\) falls back to a function of distance; the threshold test needs
  only \(L^0\) and stays available.
\end{itemize}

\subsection{Road against rail}

The specification is estimated separately for road and rail barriers.
\begin{itemize}
\item \textbf{Road:} bundled with the project and lower near-road air pollution. ÅDT on the
  shielded segment is an outcome (does traffic change at \(\tau_b\)?) and a control.
\item \textbf{Rail:} close to noise only. Similar road and rail estimates argue against air
  pollution as the channel.
\item \textbf{Air pollution as an outcome:} if a modelled NO\textsubscript{2} or particulate grid
  covers the period, its change at \(\tau_b\) by side measures what a road barrier did to
  pollution \autocite[the analogues in the lead studies are traffic metals in moss and industrial
  emissions:][]{gronqvistUnderstandingHowLow2020,hollingsworthThousandCutsCumulative2025}.
\end{itemize}

% ===========================================================================
\section{Mechanisms and other estimands}\label{sec:mechanisms}

\subsection{Mobility and the effect per year lived}\label{sec:mobility}

With assignment fixed at \(\rho_b\), moving is a \textbf{mechanism}, not a threat. On
household-years in the same stacks:
\[
M_{ft} = \sum_{k \neq -1} \mu_k\, \mathbf{1}\{t - \tau_b = k\}\, P_f + \gamma_{b,P_f,r(f)} + \eta_{b,t} + X_f'\phi + u_{ft},
\]
\(M_{ft} = 1\) if household \(f\) leaves its reference cell between \(t-1\) and \(t\), among
households still there at \(t-1\).
\begin{itemize}
\item \(\mu_k < 0\) for \(k \ge 0\): the barrier reduces moving out.
\item \(\mu_k \neq 0\) between \(\rho_b\) and \(\tau_b\): anticipation of the announcement or
  the construction work, what the early reference date guards against.
\end{itemize}

\paragraph{Per year lived (2SLS).} \(D_i^{\text{act}}\): years before the test the child
actually lived in a cell protected by any barrier, instrumented with \(D_i P_i\), with the main
specification's FE and controls.
\begin{itemize}
\item The estimate is \(\beta / \pi\), \(\pi\) the effect of assigned on realised dose: an
  average causal response among compliers \autocite{angristTwoStageLeastSquares1995}, children
  who live behind a barrier longer because their reference cell is protected.
\item Monotonicity is plausible: assignment should not reduce anyone's years behind a barrier;
  comparison children who later move behind another barrier are always-takers.
\item Exclusion binds: the reference cell must affect scores only through years behind a
  barrier. It fails if the barrier changes moving and moving changes outcomes (school,
  neighbourhood, peers). If \(\mu_k\) and the change-of-school event study are near zero after
  \(\tau_b\), we report a \textbf{LATE}; otherwise a \textbf{rescaled ITT}. Both with bounds
  allowing a direct effect of assignment \autocite{conleyPlausiblyExogenous2012}, calibrated
  from the estimated change in moving and literature effects of moving.
\item Compliers are described by household characteristics at \(\rho_b\)
  \autocite{abadieSemiparametricInstrumentalVariable2003}, presumably settled families who stay.
\item \(D_i^{\text{act}}\) comes from yearly addresses on the 100\,m grid, so it is misclassified
  near zone edges, non-classically; dropping edge cells limits the bias in \(\pi\).
\end{itemize}

\subsection{Siblings}

Household FE \(\gamma_{f(i)}\) identify the effect from siblings with different exposure, mainly
families with children tested before and after \(\tau_b\). The younger sibling always has the
higher dose, so dose is confounded with birth order \autocite{blackMoreMerrierEffect2005}:
birth order \(\times\) sex FE, identified from comparison families. We report the number of
families with siblings on both sides of \(\tau_b\). The sample leans toward large families that
stay: robustness only.

\subsection{Within child: grade 3 to grade 9}\label{sec:withinchild}

Grade-3 and grade-6 results exist from 2012 and grade-9 grades to 2022. Children born
2002--2006 have all three: grade 3 in 2012--2016, grade 6 in 2015--2019, grade 9 in 2018--2022.
For a wall built between a child's tests, the same child is observed before and after it.

\paragraph{Sample.} The stacks of walls built 2014--2021 (549 road and 364 rail records,
provisional), children born 2002--2006, assigned at \(\rho_b\) as in the main design.

\paragraph{Specification.} For child \(i\) and grade \(g \in \{3, 6, 9\}\), the intended test year
is \(t^*_{ig} = a_i + \{10, 13, 16\}\) and event time \(e_{ig} = t^*_{ig} - \tau_b\):
\[
S_{ig} = \theta\, \mathbf{1}\{e_{ig} \ge 1\}\, P_i + \alpha_i + \eta_{b,g,t} + \varepsilon_{ig}
\]
\begin{itemize}
\item \(S_{ig}\): the child's percentile among all children with the same measure in the same
  grade and year.
\item \(\alpha_i\): child FE. They absorb the side FE, the household controls and anything fixed
  about the child, including sorting into cells on levels.
\item \(\eta_{b,g,t}\): barrier \(\times\) grade \(\times\) test year FE, shared by both sides.
\item Tests at \(e_{ig} \in \{-1, 0\}\) are transition tests and are dropped, as in the main
  design. \(\theta\) is identified from children with a test before (\(e \le -2\)) and a test
  after (\(e \ge 1\)).
\item \(\theta\): the change in a child's rank after the wall, protected side against opposite
  side. It measures the immediate effect at school ages; an event-time version replaces the
  indicator with years since construction at the test.
\item \textbf{Placebo:} children with grades 3 and 6 both before the wall; their change from
  grade 3 to 6 should not differ by side.
\end{itemize}

\paragraph{What it adds and what it cannot do.}
\begin{itemize}
\item It compares each child with themselves, so it does not rely on cohorts being comparable
  across \(\tau_b\), as the main design does. Agreement of \(\theta\) and \(\delta^{\text{SR}}\)
  supports both.
\item It uses recent walls, which the main design can use only in the short window.
\item It covers school ages 10--16 only: nothing on early childhood or in utero.
\item The three measures differ (grade-3 tests cover fewer subjects; grade 9 is a merit value).
  Ranking within measure puts them on one scale, and the difference between sides nets out what
  they share, but an effect that differs by measure is not separated from an effect that grows
  with age.
\item Five birth cohorts only.
\end{itemize}

\subsection{Residents at the time of the test}

For comparison with dwelling-based studies, a repeated cross-section assigned by \textbf{where
children live in the test year}:
\[
\bar S_{b r p t} = \sum_{j} \beta_j\, \mathbf{1}\{r = j\}\, \mathbf{1}\{p = 1\}\, \mathbf{1}\{t \ge \tau_b\} + \gamma_{b r p} + \eta_{b t} + \varepsilon_{b r p t},
\]
\(b\) barrier, \(r\) distance band, \(p\) side, \(t\) test year.
\begin{itemize}
\item It mixes the direct effect with sorting; its gap to the ITT measures sorting.
\item A minimum number of test takers per cell; the test of
  \textcite{santannaDifferenceinDifferencesCompositionalChanges2025} for compositional change.
\item By side and band: the share of school-age children, enrolment and transfer rates after
  \(\tau_b\).
\end{itemize}

% ===========================================================================
\section{Threats and checks}\label{sec:threats}

\begin{center}
\footnotesize
\renewcommand{\arraystretch}{1.25}
\begin{tabular}{@{}L{0.27\textwidth}L{0.35\textwidth}L{0.32\textwidth}@{}}
\toprule
\textbf{Threat} & \textbf{Design answer} & \textbf{Evidence we show} \\
\midrule
Families sort or anticipate & Reference cell at \(\rho_b\), via the mother; ITT & Mobility event study; births by side; predicted-score placebo \\
Construction noise before \(\tau_b\) & \(e = -1, 0\) as transition; reference \(e = -2\) & \(\beta_{-1}\), \(\beta_0\) \\
Road project comes with the barrier & Opposite side of the same road; barrier \(\times\) year FE & Threshold, intensity, plateau; barrier \(\times\) band \(\times\) year FE \\
Air pollution moves with noise & Rail against road & NO\textsubscript{2} by side, if a grid exists \\
Staggered timing, bad comparisons & Stacking, own cells only, corrective weights & Callaway--Sant'Anna, Sun--Abraham, Borusyak et al. \\
Cells misclassified at the zone edge & Drop edge cells (100\,m grid) & Estimates with them included \\
Comparison cells insulated & --- (bias toward zero) & Insulation records, if they exist \\
Grading changes & Percentile within test year & Raw merit value, national tests \\
Sides feed different schools; grading by school & Assigned school \(\times\) year FE (not actual school: post-treatment) & School overlap of the two sides; actual school \(\times\) year FE as decomposition \\
\bottomrule
\end{tabular}
\end{center}

\paragraph{Threats in more detail.}
\begin{itemize}
\item \textbf{Treatment measurement.} Treatment and residence are on 100\,m SCB cells, not
  dwellings. Near a zone edge, 13--18\% of points have the opposite status from their cell. Main
  estimates drop cells straddling the edge (a donut); we also report them included.
\item \textbf{Missing construction years.} Missing for about 57\% of road and 31\% of rail walls,
  not at random: most undated road records are the batch loaded when the register started in
  2006 (built by 2006, year unknown). Dated walls are mostly recent (71\% of dated road and 54\%
  of rail records built 2015 or later) and cannot enter the full window. The full-window event
  study rests mainly on rail walls (313 of the 364 records built 2000--2005, provisional). Results are for the dated sample;
  dated and undated barriers are compared on location, length, road type and neighbourhood.
\item \textbf{Current walls only.} Demolished or replaced walls are missing; a rebuilt wall may
  carry the rebuild year.
\item \textbf{Insulation in the comparison group}: most likely a bias toward zero (\S\ref{sec:literature}).
\item \textbf{No unbuilt barriers}: opposite-side and farther-behind cells take that role.
\item \textbf{Outcome data.} The merit value is the main outcome, national tests secondary.
  Grade-3 and grade-6 results exist from 2012, which allows child FE for five cohorts
  (\S\ref{sec:withinchild}); what they measure is to be confirmed.
\item \textbf{Grading by school.} Teachers set the merit value, and the national tests are also
  marked at the school. Answered by assigned-school FE and the school overlap of the sides.
\item \textbf{Wind.} Wind direction barely changes year to year and carries air pollution like
  noise, so it cannot separate the two: heterogeneity by side (downwind/upwind) only, if at all.
\end{itemize}

\paragraph{Balance and composition checks.}
\begin{itemize}
\item \textbf{Predicted score:} regress scores on household characteristics at \(\rho_b\) in the
  comparison group, predict for everyone, and use it as the outcome. Any \(\beta_e \neq 0\)
  shows composition differing by side and cohort.
\item \textbf{Fertility and household formation:} births after \(\rho_b\) may respond to the
  barrier. Number of births per reference household and the mothers' characteristics, by
  \(P_i\) after \(\tau_b\).
\item \textbf{Placebo cohorts:} \(\beta_e\) for \(e < 0\); and \(\tau_b\) shifted back five years.
\item \textbf{Off-age testing} by \(P_i\) after \(\tau_b\); children tested on time only.
\item \textbf{Birth outcomes:} birthweight and gestational age (Medical Birth Register) for
  cohorts born after \(\tau_b\), by \(P_i\), testing the in-utero channel directly, as both lead
  studies do \emph{(access to be confirmed)}.
\end{itemize}

\paragraph{Heterogeneity, fixed in advance} so it does not read as a search: sex; parental
education, family income and immigrant background at \(\rho_b\); housing type (apartment or
house, since insulation and retrofits differ by it). In the lead literature effects concentrate
among boys and low-income families \autocite{gronqvistUnderstandingHowLow2020} and in poorer
counties \autocite{hollingsworthThousandCutsCumulative2025}.

% ===========================================================================
\section{Data and open questions}\label{sec:data}

\paragraph{Data.}
\begin{itemize}
\item \textbf{Barrier counts are provisional.} Every count of barriers, records or cells in this
  document comes from Trafikverket's register as fetched for the project (checked 2026-11-02).
  New data are discussed with Trafikverket on 25 November 2026; construction years and
  heights may change.
\item \textbf{Treatment grid} (built in the repository): 100\,m SCB cells with the
  \texttt{protected}-tier treatment, construction year, timing unknown, side unknown and the
  control band; exported to the shared folder 2026-10-22.
\item \textbf{Barriers:} Trafikverket's register; construction years for 49\% of road and 73\%
  of rail records. Of the undated, 978 road records entered NVDB's first load in 2006 (built by
  then) and 451 rail records carry BIS's first-load connection date (before April 2019).
\item \textbf{Registers (MONA)}, as stated 2026-11-18:
  \begin{itemize}
  \item grade-9 grades, 1996--2022;
  \item residences of the entire population, from 1980;
  \item grade-3 and grade-6 results, from 2012;
  \item mother links and household characteristics at \(\rho_b\) (to be confirmed).
  \end{itemize}
  Residences from 1980 never bind: \(\rho_b = \tau_b - 3 \ge 1980\) for every dated wall. The
  grade years bound the windows (decision 1).
\item \textbf{To build in the repository for this design:}
  \begin{itemize}
  \item \textbf{the stack file}: one row per cell \(\times\) barrier with the barrier and its
    kind, \(\tau_b\), \(\rho_b\), the cell's role (protected, opposite side, 600--1,500\,m band),
    its distance from the road and band, the wall it belongs to, an edge flag for the donut, and
    its group of connected corridors. The current export has one row per cell and no
    opposite-side flag;
  \item \textbf{pre-construction traffic} for \(L^0\): the traffic panel matches about 30--40\%
    of road segments; historical train counts are not yet checked.
  \end{itemize}
\end{itemize}

\paragraph{Questions for Trafikverket.} The list is kept with the barrier data, in the
repository (\url{docs/data/sweden/noise_barriers/README.md}), since the school-level line
needs the same answers. They are discussed at a call with Trafikverket on 25 November 2026
(slides in \url{output/presentations/2026-11-trafikverket/}). In short:
\begin{enumerate}
\item construction years for the undated walls, both the rail walls connected before April 2019
  (most of the short window's gain, \S\ref{sec:power}) and the road walls loaded in 2006 (the
  full window); heights for the road walls; and what the year means;
\item historical versions of the register, with removed walls;
\item the exact position of each wall: its side and distance from the road;
\item start or announcement dates of the road projects (for \(\rho_b\)), and walls planned but
  never built;
\item the register of properties considered for facade insulation: measure, type, year, cost,
  and the 100\,m grid cell;
\item noise levels calculated before and after a measure, train counts and road traffic beyond
  NVDB's history, and the qualifying rules over time.
\end{enumerate}

\paragraph{Open questions for the register data (MONA).}
\begin{enumerate}
\item Which years of individual grade-9 results (merit value, national tests) exist, and are
  children linked to their mothers? \emph{Partly answered 2026-11-18: grade-9 grades
  1996--2022. Open: which measures cover those years, and the link to mothers.}
\item From which year are residences on the 100\,m grid? \emph{Answered 2026-11-18: from 1980,
  for the entire population.}
\item Exact years of the grading reforms and the merit-value scale, for the percentile.
\item Is the grade-9 school identified for each child, and are siblings linked (assigned-school
  FE)?
\item Access to the Medical Birth Register (a separate application to Socialstyrelsen).
\item Are the grade-3 and grade-6 results from 2012 national tests or grades (grades start in
  year 6)? Children born 2002--2006 have grade 3 (2012--2016) and grade 9 (2018--2022), which
  would allow a within-child design for walls built 2012--2022.
\end{enumerate}

\paragraph{Open question for our own build.} Whether a modelled air-pollution grid covers our
period.

% ===========================================================================
\section{Pre-specification}\label{sec:prespec}

\subsection{Decisions}\label{sec:decisions}

\paragraph{Fixed already} (by this design, not open): \(\rho_b = \tau_b - 3\) until project dates
exist, robustness \(\tau_b - 1\) and \(\tau_b - 5\); WHO thresholds 53\,dB road, 54\,dB rail;
distance bands 0--100 / 100--250 / 250--600\,m; grade 9 on time as a second outcome;
anticipation measured by moves between \(\rho_b\) and \(\tau_b\); the heterogeneity set of
\S\ref{sec:threats}; the full-window event study only on barriers with a full window. Added 2026-11-18: the step at
\(e = 1\) as the prior and \(\delta^{\text{SR}}\) on the short window \(e = -K, \dots, +3\) as the
main estimate; \(e = -1, 0\) as transition cohorts; \(e = -2\) as the reference.

\paragraph{Open.} Each with a recommendation. Fill in ``Decided'' with the date.
\begin{enumerate}
\item \textbf{\(K\) and \(L\) (the windows).} Grade-9 grades cover 1996--2022, so a window
  \([\tau_b - K,\ \tau_b + \text{top}]\) fits only walls built \(1996 + K\) to
  \(2022 - \text{top}\). Records and dated protected cells per window (register as of
  2026-11-02, provisional until Trafikverket's reply):

  \smallskip
  {\footnotesize
  \begin{tabular}{@{}llrrrr@{}}
  \toprule
  window & walls built & road records & road cells & rail records & rail cells \\
  \midrule
  short, \(e = -5 \dots 3\) & 2001--2019 & 471 & 5,110 & 717 & 6,972 \\
  short, \(e = -3 \dots 3\) & 1999--2019 & 490 & 5,401 & 829 & 8,043 \\
  full, \(e = -3 \dots 17\) & 1999--2005 & 62 & 761 & 328 & 3,171 \\
  full, \(e = -4 \dots 17\) & 2000--2005 & 51 & 607 & 313 & 3,066 \\
  full, \(e = -5 \dots 17\) & 2001--2005 & 43 & 470 & 216 & 2,100 \\
  full, \(e = -4 \dots 16\) (no in utero) & 2000--2006 & 61 & 755 & 338 & 3,693 \\
  plateau, \(e = -4 \dots 20\) & 2000--2002 & 25 & 391 & 137 & 1,538 \\
  \bottomrule
  \end{tabular}}
  \smallskip

  Recommended:
  \begin{itemize}
  \item \textbf{short window \(K = 5\)} (\(e = -5, \dots, 3\)): three placebo cohorts beyond the
    reference, for few walls lost;
  \item \textbf{full window \(K = 4\), \(L = 1\)} (\(e = -4, \dots, 17\)): two placebo cohorts
    and the in-utero step. \(K = 3\) leaves a single placebo cohort (\(e = -3\)), and \(K = 5\)
    loses a third of the rail walls;
  \item \textbf{no plateau as a target}: \(L = 4\) keeps 25 road records. The full window is
    mostly a rail result unless Trafikverket recovers years for the road walls loaded in 2006:
    4 usable road stacks today, about 46 with the missing years, at most 164
    (\S\ref{sec:power}).
  \end{itemize}
  \textbf{Decided:}
\item \textbf{The percentile.} Recommended: merit value ranked among all children in Sweden
  tested in the same year, 0--100, ties at the mid-rank. Alternative: rank within birth cohort
  (mixes in off-age children).\\ \textbf{Decided:}
\item \textbf{Donut width.} Recommended: drop every cell whose 100\,m square crosses the edge
  of a protected or comparison area. Alternative: cells whose centroid is within 50\,m of an
  edge. Report both.\\ \textbf{Decided:}
\item \textbf{Clusters.} Recommended: groups of connected corridors, with WCR and RI.
  Alternative: kommun.\\ \textbf{Decided:}
\item \textbf{Barriers walled on both sides} (about 40\% of those built by 2007). Recommended:
  they enter only through the 600--1,500\,m band; the opposite-side estimate uses one-sided
  barriers. Report the share of protected children this moves. If decision~\ref{dec:comparison}
  pools the band, these barriers enter the main estimate through it.\\ \textbf{Decided:}
\item\label{dec:profile} \textbf{If the distance profile has not faded by 600\,m.} The strategy
  so far moves the band out, and lets the profile guide the protected zone. Both choose a sample
  after seeing outcomes. Recommended: keep the protected zone as the \texttt{protected} tier
  (600\,m); if the profile's 600--1,500\,m bins are jointly non-zero (\(p < 0.10\)), drop
  comparison group 2 and report it, rather than moving the band. If
  decision~\ref{dec:comparison} pooled it, the main estimate falls back to the opposite side
  alone.\\ \textbf{Decided:}
\item \textbf{The mobility event study's reference period.} Households are defined by living in
  the reference cell at \(\rho_b = \tau_b - 3\), so their move hazard starts at \(k = -2\): the
  reference \(k = -1\) lies inside the anticipation window, and the side FE absorb any level
  difference from \(\rho_b\) on. As written, the equation cannot show anticipation.
  Recommended: add a pre-period from households living in the same cells at an earlier date
  (\(\tau_b - 8\), followed to \(\tau_b - 3\)), so the hazard between \(\rho_b\) and \(\tau_b\)
  is compared, by side, with the same cells' hazard before \(\rho_b\).\\ \textbf{Decided:}
\item \textbf{The within-child design} (\S\ref{sec:withinchild}). Recommended: include it as a
  secondary block (rule F1), run once MONA confirms what the grade-3 and grade-6 results are
  (question 6 of \S\ref{sec:data}). Alternative: leave it for a later pre-specification.\\
  \textbf{Decided:}
\item\label{dec:comparison} \textbf{The comparison in the main estimate.} The opposite side alone
  holds few children (\S\ref{sec:power}). Recommended: the opposite side and the
  600--1,500\,m band pooled in each stack (comparison groups 1 and 2), with the FE of
  \S\ref{sec:spec}; the opposite side alone (barrier \(\times\) band \(\times\) event time FE)
  as the robustness check, and rule B3 comparing the two. Short-run MDE, opposite side against
  both: road 0.46 against 0.22 SD, rail 0.25 against 0.15 SD. The cost: the band shares the
  works and the traffic changes less closely, so a project effect that fades with distance can
  load onto \(\beta_e\); B3 measures it. Alternative: the opposite side alone, as drafted.\\
  \textbf{Decided:}
\item\label{dec:primary} \textbf{One primary estimate.} Rules A2 and B2 estimate road and rail
  separately but do not say which carries the main claim. Recommended: rail is primary (rule A2
  at \(p < 0.05\), no correction) and road secondary, labelled as such. Rail is close to noise
  only and more precise (short-run MDE 0.15 against 0.22 SD with both comparisons). Alternative:
  both primary, Holm over the two.\\ \textbf{Decided:}
\end{enumerate}

\subsection{Rules}\label{sec:rules}

Construction rules (G) load no outcome and run in the repository. The others run on MONA.
\begin{itemize}[leftmargin=4.2em]
\prule{G1} Only barriers with a construction year enter. Before estimation, dated and undated
  barriers are compared on location, length, road type and neighbourhood; the table is
  recorded in \S\ref{sec:outcomes}.
\prule{G2} An opposite-side cell is in the same distance band from the same road segment, on the
  side the barrier does not shield, and protected by no barrier within the stack's window. A
  comparison cell protected by another barrier within the window is dropped from that stack.
\prule{G3} Cells in the donut (decision 3) carry the edge flag and are left out of the main
  estimate.
\prule{G4} Before anything is sent to MONA, record: barriers per kind with a full window (given
  decision 1), protected and comparison cells per role and band, corridor groups and their
  sizes, and barriers walled on both sides.
\prule{G5} \(L^0\) is computed only where pre-construction traffic exists; the threshold and
  intensity tests run on those barriers, and their count is recorded.
\prule{A1} \emph{Placebo.} If \(\beta_e\), \(e \le -3\), are jointly non-zero at \(p < 0.10\)
  (WCR by corridor group), nothing after is interpreted; label: ``pre-trend''. \(\beta_{-1}\) and
  \(\beta_0\) are reported, not tested as placebos; if negative at \(p < 0.10\), label:
  ``construction dip''.
\prule{A2} \emph{Main claim.} ``Living behind a barrier changes the grade-9 percentile at once''
  is claimed only if \(\delta^{\text{SR}}\) (short window; for the primary kind of
  decision~\ref{dec:primary}) has RI \(p < 0.05\) (permuting
  \(\tau_b\) across corridor groups) and A1 passes. Otherwise: ``no short-run effect detected''.
\prule{A3} \emph{Shape} (full window), only if A1 passes there. Each label stands on its own:
  ``step'' if \(\delta\) has \(p < 0.05\); ``cumulative'' if \(\beta > 0\) at \(p < 0.05\);
  the plateau is not tested with \(L = 1\) (decision 1); if \(L \ge 2\) is chosen, ``plateau''
  if \(\beta_e\) is flat after \(e = A + 1\) (slope \(p \ge 0.10\)), otherwise ``no plateau''
  (points to the project).
\prule{B1} \emph{Noise, not the project.} Claimed only if \(\delta^{\text{SR}}\) above the
  threshold has \(p < 0.05\), below it \(p \ge 0.10\), and their difference \(p < 0.05\).
\prule{B2} Road and rail are estimated separately and both reported; a difference is a finding
  only at \(p < 0.05\).
\prule{B3} \emph{Comparison gap} (if decision~\ref{dec:comparison} pools the band).
  \(\delta^{\text{SR}}\) is also estimated against the opposite side alone and against the band
  alone. If they differ at \(p < 0.05\), the main estimate is labelled ``depends on the
  comparison'', and the opposite-side estimate, which nets out the project, is reported next to
  it.
\prule{C1} \emph{Anticipation} is found if moves out of protected reference cells between
  \(\rho_b\) and \(\tau_b\) differ from the opposite side at \(p < 0.05\). Then the
  \(\tau_b - 5\) estimate is reported next to the main one, labelled ``may include
  anticipation''.
\prule{D1} The five heterogeneity cuts are the only ones claimed, with Holm over the five. Wind
  by side, if run, is labelled exploratory.
\prule{E1} Grade 9 on time is a secondary outcome with its own rule A2, never combined with the
  main outcome.
\prule{F1} \emph{Within child} (if decision 8 includes it). ``A child's rank changes after the
  wall'' is claimed only if \(\theta\) has RI \(p < 0.05\) (permuting \(\tau_b\) across corridor
  groups) and the grade 3--6 placebo shows no difference by side (\(p \ge 0.10\)). Labelled
  ``school ages, within child''; secondary, never combined with A2. If its sign differs from
  \(\delta^{\text{SR}}\), both are reported as conflicting.
\end{itemize}

\subsection{Outcomes}\label{sec:outcomes}

Nothing has run.

\subsection{Change log}\label{sec:changelog}

\begin{itemize}
\item \textbf{2026-11-18, draft.} Merged from \texttt{empirical\_strategy.tex} (the design), the
  slides revised 2026-11-04 (distance bands, grade 9 on time, anticipation by moves, barriers
  walled on both sides) and the first Markdown draft of this pre-specification (decisions and
  rules). Decision 6 is new: it replaces moving the band, and the profile guiding the protected
  zone, with a rule fixed in advance.
\item \textbf{2026-11-18, draft.} The prior is an immediate effect of noise. Added: the step at
  \(e = 1\) (\(\delta\)) next to the ramp (\(\beta\)); \(e = -1, 0\) as transition cohorts
  (construction works; the wall may finish after the test); the reference cohort moves from
  \(e = -1\) to \(e = -2\); a short balanced window \(e = -K, \dots, +3\) with the short-run effect
  \(\delta^{\text{SR}}\) as the main estimate (rules A1--A3 and B1 rewritten); decision 7 on the
  mobility event study's reference period. No outcome has been loaded.
\item \textbf{2026-11-18, data.} \S\ref{sec:data} records the register years (grade 9
  1996--2022, residences from 1980, grades 3 and 6 from 2012), marks the barrier counts as
  provisional until Trafikverket's reply, and moves the Trafikverket questions to the
  repository's barrier data page (summarised here).
\item \textbf{2026-11-18, windows and within-child design.} Decision 1 now carries the windows the
  grade years allow, with counts, and recommends \(K = 5\) (short), \(K = 4\), \(L = 1\) (full)
  and no plateau target; the full window is walls built 2000--2005, not ``by 2007''. Rule A3's
  plateau label follows decision 1. New: the within-child design (\S\ref{sec:withinchild}),
  decision 8 and rule F1. No outcome has been loaded.
\item \textbf{2026-11-19, power.} New \S\ref{sec:power}: group sizes and MDEs from open
  population data, today and with every construction year. From it: the opposite side's share
  is stated in children (27\% road, 45\% rail), not cells (``about 60\%''); new decisions 9
  (the band pooled with the opposite side in the main estimate) and 10 (rail as the primary
  estimate), with rule B3 and rule A2 pointing to decision 10; decisions 1, 5 and 6 note what
  follows; the Trafikverket questions ask for the rail years as well as the road years. No
  outcome has been loaded.
\end{itemize}

\printbibliography[title={References}]

\end{document}
